\chapter{CONCLUSÃO}
\label{chapter-7}
O Prazo Claro materializou uma aplicação web acadêmica de contagem processual com resultados explicáveis e limites identificados. A solução separou a entrada e a seleção de dados do algoritmo determinístico, preservou versões e registrou a trilha de cálculo. A evolução de um recorte civil para perfis criminais e trabalhistas ocorreu mediante condições próprias, sem apresentar um método genérico como válido para qualquer prazo.

Os objetivos de especificação, implementação e verificação técnica foram alcançados no recorte documentado. Conta, histórico e lembrete opcional demonstram integração com serviços externos, enquanto as evidências diferenciam execução local, resposta simulada, uso hospedado e recebimento real. O conjunto de testes aprovado sustenta o comportamento verificado, mas não substitui revisão jurídica ou avaliação ampla com usuários.

A contribuição de engenharia está na combinação de reprodutibilidade, comunicação de incompletude e rastreabilidade entre requisitos, código e evidências. A documentação permite instalar, demonstrar e continuar o artefato sem confundir recursos planejados com recursos existentes. O resultado permanece demonstrativo: calendário completo, aprovação jurídica, fechamento do ensaio externo e preparação institucional para operação ampliada constituem etapas posteriores claramente delimitadas.
